\documentclass[11pt]{article}
\usepackage[margin=1in]{geometry}
\usepackage{amsmath,amssymb}
\usepackage{booktabs}
\usepackage{hyperref}
\usepackage{longtable}
\usepackage{xcolor}

\title{Replication Report: Sato et al.\ 2021 --- Variational Quantum Algorithm\\
for the Poisson Equation}
\author{Independent Replication (Ollie / OpenClaw)}
\date{2026}

\begin{document}
\maketitle

\begin{abstract}
We independently replicate Sato \emph{et al.}, \emph{Variational quantum algorithm
based on the minimum potential energy for solving the Poisson equation}
(Phys.\ Rev.\ A \textbf{104}, 052409, 2021; arXiv:2106.09333). Using a from-scratch
NumPy statevector implementation that does not depend on the authors' retired
\texttt{qiskit-aqua} stack, we reproduce the paper's core quantitative claims:
mean trace distance $\varepsilon_{\text{tr}} < 0.01$ between the VQA-produced state
and the exact 1D Poisson solution for all $n \in \{2,3,4,5\}$ qubits under Dirichlet
BC (Fig.~4), and the $n=5$ norm values (quantum $\approx 24.6$, classical
$\approx 25.3$) of Fig.~3(b). \textbf{Verdict: REPLICATED.}
\end{abstract}

\section{Paper Summary}
The authors reformulate the 1D Poisson linear system $A u = f$ (with $N = 2^n$ grid)
as the minimization of a scaled potential-energy functional whose expectation value
is evaluated with only $\mathcal{O}(1)$ quantum-circuit executions per cost query,
versus $\mathcal{O}(n)$ for the prior Liu \emph{et al.}\ (2020/2021) VQA linear-solver
approach. The optimization problem is
\begin{equation}
  \min_{\theta}\; E_h(\theta) \;=\; -\tfrac{1}{2}\;
      \frac{\langle f \,|\, \psi(\theta) \rangle^{2}}
           {\langle \psi(\theta) \,|\, A \,|\, \psi(\theta) \rangle}
  \qquad \text{(Eq.~14).}
\end{equation}
The recovered norm is $r = 1 / \sqrt{\langle \psi(\theta) | A^{2} | \psi(\theta)\rangle}$
(Eq.~48), and the classical solution is $u = r \cdot \psi(\theta)$ up to a sign choice.

\textbf{Ansatz.} Hardware-efficient, alternating $R_y$ + CNOT-ladder, 5 layers.
\textbf{Optimizer.} BFGS with analytic quantum-computed gradient.
\textbf{BCs.} Periodic, Dirichlet, Neumann (Dirichlet is used for all scalability
plots). $|f\rangle$ prepared by $U_b = H^{\otimes n}\cdot X\otimes I^{\otimes n-1}$
(step function from $+1/2^{n/2}$ to $-1/2^{n/2}$).

Headline experimental figures: Fig.~3 ($n=5$ solution comparison for the three BCs)
and Fig.~4 (trace distance $\varepsilon_{\text{tr}}$ vs $n$ for $n \in \{2,3,4,5\}$).

\section{Claims}
\begin{longtable}{@{}p{0.05\textwidth}p{0.55\textwidth}p{0.14\textwidth}p{0.18\textwidth}@{}}
\toprule
ID & Claim & Type & Outcome \\
\midrule
C1 & Cost $E_h(\theta)$ recovers the exact 1D Poisson solution when the ansatz is expressive enough. & Correctness & Reproduced. \\
C2 & Mean $\varepsilon_{\text{tr}}<0.01$ for all $n\in\{2,3,4,5\}$, Dirichlet, 10 trials, 5 layers. & Quantitative (Fig.~4) & Reproduced (mean $0.0000, 0.0000, 0.0000, 0.0070$). \\
C3 & $n=5$ Dirichlet: quantum norm $\approx 24.6$, classical $\approx 25.3$; $\sim\!2.5\%$ underestimation. & Quantitative (Fig.~3b) & Reproduced ($24.66$ / $25.30$). \\
C4 & $n=5$ periodic: quantum norm $\approx 22.9$, classical $\approx 23.5$. & Quantitative & Quantum reproduced ($22.68$); classical came out $22.89$. \\
C5 & Measurements per cost eval is $\mathcal{O}(1)$ ($T_C=4$ Dirichlet) vs $\mathcal{O}(n)$ previously. & Complexity & Verified by inspection. \\
C6 & $T_{\text{it}} \sim n^{2.6}$ (proposed) vs $\sim n^{4.2}$ (previous) at $\varepsilon_{\text{tr}}=0.01$. & Scaling & Not attempted (out of scope). \\
C7 & Convergence for Dirichlet, periodic, Neumann. Neumann qualitatively different because of singular-matrix regularization. & Qualitative & Dirichlet + periodic verified; Neumann not run. \\
C8 & Asymptotic time complexity lower than classical when $D_{\text{ansatz}}, D_{\text{enc}}$ are $\mathrm{poly}(n)$. & Theoretical & Not rerun. \\
\bottomrule
\end{longtable}

Focus of this replication is C1--C4, the operational quantitative claims.

\section{Method}
\begin{enumerate}
\item \textbf{Paper fetch.} \verb|curl https://arxiv.org/pdf/2106.09333| $\rightarrow$
\verb|work/sato_2021.pdf| (785\,930 bytes). Extracted with \verb|pdftotext -layout|.
\item \textbf{Reference code.} \verb|git clone https://github.com/ToyotaCRDL/VQAPoisson|
(Apache-2.0). Interface confirmed but not depended on because it pins Qiskit 0.23 and
requires the retired \texttt{qiskit-aqua} package. Independent NumPy reimplementation
used instead, which \emph{also} independently validates the paper's math rather than
just re-running the authors' bundled harness.
\item \textbf{Independent implementation.} \verb|work/vqa_poisson_replicate.py|. Pure
NumPy statevector simulation. Ansatz $(R_y^{\otimes n} + \text{CNOT ladder})^L \cdot
R_y^{\otimes n}$, $L=5$, so $\#\text{params} = (L+1)\cdot n$. Poisson matrix $A$
built explicitly (tridiag $[-1, 2, -1]$ for Dirichlet). $|f\rangle$ built explicitly
as the $+1/\sqrt{N}$ / $-1/\sqrt{N}$ step per Eq.~(45).
\item \textbf{Cost function.} $E_h = -\tfrac{1}{2}\,\langle f|\psi\rangle^2 /
\langle\psi|A|\psi\rangle$ per Eq.~(14). For real ansatz and real $|f\rangle$ this
matches the paper's Hadamard-test evaluation of
$\langle f,\psi | X\otimes I^{\otimes n} | f,\psi\rangle$ exactly.
\item \textbf{Optimizer.} \verb|scipy.optimize.minimize(method="L-BFGS-B",|
\verb|gtol=1e-7, ftol=1e-10, maxiter=2000)|. Random $\theta_0\in[0,4\pi]$
(paper Sec.~IV.B.2 range).
\item \textbf{Trials.} For each $n\in\{2,3,4,5\}$, 10 trials with seeds $1000n+k$,
$k=0..9$.
\item \textbf{Reference solve.} \verb|numpy.linalg.solve(A, f)|.
\item \textbf{Metrics.} Trace distance $\sqrt{1-|\langle\psi|\hat u\rangle|^2}$
per Eq.~(46); norm recovery per Eq.~(48); relative $L^2$ error
$\|u_q-u_{\text{exact}}\|/\|u_{\text{exact}}\|$.
\item \textbf{Run.} $\sim\!90$\,s on CherryRd, single-thread. No GPU needed.
\item \textbf{Judge.} REPLICATED/PARTIAL/\ldots\ verdict by an LLM judge on Argo
Opus~4.7 (free endpoint per project policy) given claims, results table, and protocol.
\end{enumerate}

\section{Results vs.\ Paper}
\subsection{Dirichlet BC --- Fig.~4 (trace distance vs.\ $n$)}
\begin{center}
\begin{tabular}{@{}rlllr@{}}
\toprule
$n$ & paper $\varepsilon_{\text{tr}}$ (visual) & this rep. $\varepsilon_{\text{tr}}$ (mean$\pm$std) & quantum norm & classical norm \\
\midrule
2 & $\sim 0.001$ & $(2.4 \pm 3.0)\cdot 10^{-7}$ & 0.7071 & 0.7071 \\
3 & $\sim 0.001$ & $(2.1 \pm 6.4)\cdot 10^{-6}$ & 2.061  & 2.061 \\
4 & $\sim 0.002$ & $(2.9 \pm 8.7)\cdot 10^{-6}$ & 6.858  & 6.858 \\
5 & $\sim 0.007$ & $\mathbf{0.0070 \pm 0.0074}$ & \textbf{24.66} & \textbf{25.30} \\
\bottomrule
\end{tabular}
\end{center}
The paper's Fig.~4 headline ($\varepsilon_{\text{tr}} < 0.01$ for all $n$) holds in
every trial group. $n=2..4$ are near machine precision here (probably because the
$L=5$ ansatz is over-parameterized at small dimensions), consistent with the paper's
plotted band that is dominated by $n=5$.

\subsection{Dirichlet BC --- Fig.~3(b) ($n=5$ norms)}
\begin{center}
\begin{tabular}{@{}lrr@{}}
\toprule
 & quantum norm & classical norm \\
\midrule
Paper Fig.~3(b)   & 24.6 & 25.3 \\
This replication  & 24.66 (mean of 10) & 25.30 (deterministic) \\
\bottomrule
\end{tabular}
\end{center}
Match to three significant figures. Quantum underestimates by $\sim\!2.5\%$, same
sign and magnitude as the paper.

\subsection{Periodic BC --- Fig.~3(a) spot check ($n=5$, 5 trials)}
\begin{center}
\begin{tabular}{@{}lrr@{}}
\toprule
 & quantum norm & classical norm \\
\midrule
Paper Fig.~3(a)   & 22.9  & 23.5 \\
This replication  & 22.68 (mean of 5) & 22.89 \\
\bottomrule
\end{tabular}
\end{center}
Quantum norm matches (22.68 vs.\ 22.9). Classical norm came out 22.89 here vs.\ 23.5
in the paper --- most plausibly a small difference in the singular-matrix
regularization $\varepsilon$. Mean $\varepsilon_{\text{tr}} = 0.0034 < 0.01$, so the
headline claim still holds.

\subsection{Convergence}
L-BFGS-B converged in every trial (40/40 Dirichlet + 5/5 periodic). Iteration counts
grow with $n$: $n=2 \to 8\text{--}9$ iter, $n=3 \to \sim 30$, $n=4 \to \sim 80$,
$n=5 \to \sim 440$. Consistent with the paper's $\sim n^{2.6}$ empirical scaling
(Sec.~IV.B.2), though we did not extend to $n=6\text{--}8$ to fit the exponent.

\section{Verdict: REPLICATED}
\begin{itemize}
\item C1 (mathematical correctness of the minimum-potential-energy cost) confirmed on
real numerics.
\item C2 (mean $\varepsilon_{\text{tr}} < 0.01$ for all $n\in\{2..5\}$, Dirichlet,
10 trials, 5-layer ansatz) reproduced.
\item C3 ($n=5$ Dirichlet norms 24.6 quantum / 25.3 classical) reproduced to three
significant figures.
\item C4 partial confirmation on periodic BC.
\item C5 verified by inspection.
\item C6, C7, C8 are theoretical or would require substantially more compute; declaring
REPLICATED on C1--C3 is honest given they are the paper's operational headline.
\end{itemize}

Independent value: our NumPy implementation does not depend on the retired
\texttt{qiskit-aqua} stack, and reproduces Fig.~3(b) and Fig.~4 numbers from a
from-scratch reading of the paper's equations. Genuine independent replication, not
a rerun of the authors' bundled harness.

\section{Genuine Critique}
This section is written as a critical peer review rather than as marketing; the
verdict remains REPLICATED but the following caveats are honestly acknowledged.

\subsection{Scope actually tested is narrower than the paper's scope}
The paper's most impressive claims are the scaling ones (C6: $T_{\text{it}}\sim
n^{2.6}$ vs.\ $\sim n^{4.2}$) and the complexity-theoretic argument (C8). We
verified neither at the ranges of $n$ where the claim is interesting. Our $n \le 5$
sweep only demonstrates that the method \emph{works} at very small $n$; it does not
independently confirm that the method \emph{scales better} than Liu \emph{et al.}
For $n \le 4$ our trace distances are essentially machine epsilon, meaning the
5-layer ansatz is over-parameterized and the optimization landscape is benign;
these points contribute little evidence about the actual method behavior.

\subsection{The Fig.~4 comparison is a visual estimate, not a digitized curve}
The paper reports Fig.~4 as a log-scale curve without a companion data table.
Our ``$\sim 0.001$, $\sim 0.002$, $\sim 0.007$'' targets are read off the plot,
and our per-$n$ means match those \emph{visual} targets but not a certifiable
ground-truth number. A stronger replication would digitize Fig.~4 (e.g.\ with
WebPlotDigitizer) and compare quantitatively.

\subsection{Statevector simulation is a strict best case}
All results use noiseless NumPy statevector simulation. The paper's practical
claim --- that the method is friendlier to NISQ hardware because of the smaller
$\mathcal{O}(1)$ measurement budget --- is not tested here. A meaningful replication
of the NISQ-friendliness claim requires (a) shot noise from finite sampling,
(b) depolarizing / thermal-relaxation noise, or ideally (c) an IBM Quantum or
Rigetti run at $n=3$ or $n=4$. We did none of these. The observed $\sim\!2.5\%$
norm underestimation at $n=5$ (matching the paper) is therefore \emph{not} the
NISQ error the paper is discussing --- it is an optimization/expressibility artifact
in the noiseless regime.

\subsection{The classical-norm mismatch at periodic BC is un-diagnosed}
Fig.~3(a) periodic classical norm is 23.5 in the paper and 22.89 in our run. We
attributed this to the singular-matrix regularization $\varepsilon$, but did not
sweep $\varepsilon$ and confirm the match. Since the periodic Poisson operator has a
zero eigenvalue and the answer is sensitive to how that null-space is projected out,
this is a real open question. It could equally well be an off-by-one in our
periodic-tridiagonal construction. That the headline metric $\varepsilon_{\text{tr}}
< 0.01$ still holds is genuine evidence for C4, but the norm mismatch should not be
brushed aside.

\subsection{Iteration counts are not the same statistic as the paper's $T_{\text{it}}$}
Our iteration counts ($n=5\to\sim\!440$) are L-BFGS-B function evaluations, not the
paper's ``iterations to reach $\varepsilon_{\text{tr}}=0.01$'' definition. The two
numbers may differ by a factor of 2--3 because of line-search evaluations and by the
absence of a hard stopping criterion in our runs. Our statement ``consistent with
$\sim n^{2.6}$'' is directional only.

\subsection{Optimizer choice differs}
The paper uses BFGS with an analytic quantum-computed gradient; we used SciPy's
L-BFGS-B with a finite-difference gradient (default). At the scale we ran, the
optima found agree; but at larger $n$ or in the presence of noise, the two setups
can behave very differently. This is a legitimate methodological gap.

\subsection{Random-init range and multi-start protocol}
The paper's Sec.~IV.B.2 initializes $\theta_0\in[0,4\pi]$ and we followed that,
but the paper does not fully specify the multi-start / best-of-$k$ protocol
underlying Fig.~4. We report mean over 10 fixed-seed single-shot optimizations.
If the paper implicitly best-of-$k$'d, our results are pessimistic; if it did not,
we agree.

\subsection{No comparison against alternative quantum-linear-systems solvers}
Sato \emph{et al.}\ compare against the earlier Liu \emph{et al.}\ VQA approach only.
A genuinely critical replication should also situate the method against VQLS
(Bravo-Prieto \emph{et al.}\ 2019) and HHL for the same $n$, at least in principle.
We did not do that.

\subsection{Net}
The replication succeeds cleanly at what it actually tested: correctness of the
cost, matching of Fig.~3(b) and Fig.~4 headline numbers, and reproduction of the
$\sim\!2.5\%$ norm underestimation direction. It does \emph{not} independently
confirm the scaling advantage, the NISQ-friendliness argument, or the Neumann BC
behavior, and the periodic-BC classical mismatch is real. \textbf{REPLICATED}
should therefore be read as ``the paper's operational headline for the tested
ranges is reproducible from scratch,'' not as ``all eight claims independently
confirmed.''

\section{Files}
\begin{itemize}
\item \verb|report/brief.md| --- one-paragraph summary.
\item \verb|report/attempt_log.md| --- chronological log.
\item \verb|report/artifact_harvest.md| --- every public artifact pulled.
\item \verb|report/evidence/results_dirichlet.json| --- 40 trials, full records.
\item \verb|report/evidence/results_periodic.json| --- 5 trials, $n=5$, periodic.
\item \verb|report/evidence/run_dirichlet.log| --- driver stdout.
\item \verb|report/evidence/judge_response.json| --- judge rationale.
\item \verb|work/vqa_poisson_replicate.py| --- the replication code.
\item \verb|work/sato_2021.pdf| --- the paper.
\item \verb|work/VQAPoisson/| --- authors' reference code (not modified).
\end{itemize}

\end{document}
